%%%%%%%%%%%%%%%%%%%%%%%%%%%%%%%%%%%%%%%%%%%%%%%%%%%%%%%%%%%%%%%%%%%%%%
\section{Contexte et objectif de la phase}
%%%%%%%%%%%%%%%%%%%%%%%%%%%%%%%%%%%%%%%%%%%%%%%%%%%%%%%%%%%%%%%%%%%%%%

Le projet ALORA (\textit{Automatic Learning for Oral Resource-poor African
Languages}) vise à développer des méthodes d'intelligence artificielle
adaptées aux langues africaines peu dotées, avec le soussou comme
premier cas d'étude. L'enjeu principal est la construction d'un corpus
oral de qualité, sans lequel la modélisation du langage et de la parole
reste limitée.

Dans la plupart des langues peu dotées, l'absence de données annotées
empêche le développement de systèmes performants de reconnaissance
vocale, de transcription ou de traduction. Cette première phase a donc
pour objectif de poser les fondations techniques et scientifiques du
projet : collecter des données orales, les valider, les transcrire,
les segmenter, les annoter et les organiser dans une structure
réutilisable.

Le livrable attendu de cette phase n'est pas un modèle final, mais une
base de travail solide, reproductible et exploitable pour les étapes
suivantes du projet : corpus initial, métadonnées normalisées, pipeline
de validation et protocole de préparation des données. Cette note
synthétique vise à servir de base de discussion avec l'encadrant avant
le début de l'implémentation. Elle repose sur une architecture
modulaire, un corpus structuré et une méthodologie incrémentale, afin
d'obtenir un premier livrable exploitable pour les expérimentations
future.

\begin{important}
La phase actuelle ne vise pas à développer un modèle de reconnaissance
vocale ou de traduction, mais à construire une base de données orale
structurée, documentée et validée pour les travaux futurs.
\end{important}